\chapter{不确定性分析}

\section{局部协方差近似}
在收敛点附近，窗口优化可线性化为
\[
\Delta\bm X\approx (\bm J^\T\bm W\bm J)^{-1}\bm J^\T\bm W\,\Delta\bm z.
\]
因此位置协方差可近似为
\[
\bm\Sigma_X\approx (\bm J^\T\bm W\bm J)^{-1}.
\]
提取窗口末端 $2\times2$ 子块即可得到当前时刻位置方差估计。

\section{二维置信椭圆}
设当前时刻协方差为 $\bm\Sigma_t$，则置信区域
\[
(\bm p-\hat{\bm p})^\T \bm\Sigma_t^{-1}(\bm p-\hat{\bm p})\le \chi^2_{2,p}.
\]
对 $p=0.95$，椭圆主轴长度由特征值决定，可用于可视化定位可信度。

\section{不确定性来源分解}
误差主要来自三部分：
\begin{enumerate}
  \item RSSI 测量噪声；
  \item 路径损耗参数 $P_0,n$ 标定误差；
  \item 停车场几何遮挡导致的系统偏差。
\end{enumerate}
其中第二、三项会引入模型偏差，难以通过简单滤波完全消除。

\section{覆盖几何指标}
定义局部几何指标
\[
G_t=\operatorname{trace}\left((\bm J_t^\T\bm J_t)^{-1}\right),
\]
$G_t$ 越大说明该时刻信标几何越差。实验发现，车道尽头和立柱密集区域通常对应较大 $G_t$，与误差峰值位置一致。

\section{Bootstrap 经验区间}
为避免高斯假设过强，采用残差重采样进行 Bootstrap，生成位置估计经验分布并给出 95\% 经验区间。该方法在遮挡场景下更贴近实际误差边界。

\section{在线可信度输出}
工程上建议实时输出以下可信度字段：
\begin{itemize}
  \item 椭圆长短轴长度；
  \item 当前有效信标数量；
  \item 残差均值与最大残差；
  \item 是否触发异常保护。
\end{itemize}
这样调度端可区分“高可信定位”与“低可信定位”状态。
